\documentclass{article}
\usepackage[margin=2cm]{geometry}
\usepackage{amsmath}
\usepackage{graphicx}
\usepackage{booktabs}
\usepackage[utf8]{inputenc}
\usepackage{ragged2e}
\setlength{\emergencystretch}{3em}
\begin{document}
\sloppy

\noindent\hfill\includegraphics[height=1.8em]{assets/figure_header_logo.png}\par\noindent\hrule\vspace{4pt}

0.1 Quantile direct effect on the always-takers
We prove that


\begin{equation}
\begin{array} { r } { \theta _ { 1 } ^ { a } ( q ) = F _ { Y _ { 1 } ( 1, 1 ) | \tau = a } ^ { - 1 } ( q ) - F _ { Y _ { 1 } ( 0, 1 ) | \tau = a } ^ { - 1 } ( q ), } \\ { = F _ { Q _ { 1 1 } ( Y _ { 0 } ) | D = 0, M = 1 } ^ { - 1 } ( q ) - F _ { Y _ { 1 } | D = 0, M = 1 } ^ { - 1 } ( q ). } \end{array}
\end{equation}


This requires showing that


\begin{equation}
\begin{array} { c } { F _ { Y _ { 1 } ( 1, 1 ) | \tau = a } ( y ) = F _ { Q _ { 1 1 } ( Y _ { 0 } ) | D = 0, M = 1 } ( y ) ~ \mathrm { a n d } \qquad ( 1 ) } \\ { F _ { Y _ { 1 } ( 0, 1 ) | \tau = a } ( y ) = F _ { Y _ { 1 } | D = 0, M = 1 } ( y ). \qquad }\end{array}
\end{equation}


Under Assumptions 7 and 8.


\begin{equation}
\begin{array} { r } { F _ { Y _ { t } | D = 0, M = 1 } ( y ) = E [ 1 \{ Y _ { t } \leq y \} | D = 0, M = 1 ] } \\ { \overset { A 7, A 8 } { = } E [ 1 \{ Y _ { t } ( 0, 1 ) \leq y \} | \tau = a ] } \\ { = F _ { Y _ { t } ( 0, 1 ) | \tau = a }, ( y ). } \end{array}
\end{equation}


which provesO. From O, we have


FQ(Yo)ID=0,M=1(y) = FY(1,1)ID=0.,M(0)=1(y) = E[1\{Yi(1,1)   y\}|D = 0, M(0) = 1].


Under Assumption 7 and 8,


E[1\{Y(1, 1)   y\}]D = 0, M(0) = 1] 44* E[1\{Y(1, 1)   y\}T = a]
(4)
= FY(1,1)=a(y),


which proves

\end{document}
